% =====================================================================
% LUCRAREA MULTILINGVA, un singur fisier, de dat copy-paste in Overleaf.
% Template-ul ACL cu textul nostru in el, de la abstract la Limitations.
% Scrisa pe 19 august 2026. Tinta: ARR 12 octombrie 2026, NAACL 2027.
%
% ASTA NU E paperul romanesc. Acela e Paper/Articol_Complexity/paper_full.tex,
% alt subiect, alt set de date, si nu se atinge de aici.
%
% Numerotarea: 1 Introduction, 2 Related Work, 3 Approach, 4 Results,
% 5 Conclusions, apoi Limitations nenumerotata. Fara Acknowledgments si
% fara apendice la review.
%
% ORDINEA IN CARE A FOST SCRISA, si nu e ordinea din fisier: Results plus
% cele trei tabele ale lui, apoi Limitations, apoi Approach, apoi Related
% Work si Conclusions, apoi Introduction si abstractul. Sectiunile complet
% determinate de masuratori s-au scris primele, ca sa nu se aleaga nicio
% cifra ca sa sune bine intr-o fraza deja scrisa.
%
% ABATERILE DE LA TEMPLATE:
%
%   1. Corpul lucrarii. Abstractul template-ului e inlocuit cu al nostru,
%      \section{Introduction} cu Sectiunile 1-5, si \section*{Limitations}
%      cu a noastra. Sectiunile de manual ale template-ului (Engines,
%      Preamble, Document Body, BibTeX Files) sunt SCOASE, cu tabelele si
%      figurile lor de exemplu. Deci nu e nevoie de pachetul mwe, si dispar
%      \nocite{Ando2005,...} si citarile la Gusfield:97.
%
%   2. Bibliografia. Blocul filecontents scrie custom.bib cu cele 7 intrari
%      ale template-ului PLUS cele 19 ale lucrarii. Cele 7 raman necitate,
%      dar inofensive: fisierul merge lipit si intr-un proiect ACL care are
%      deja un custom.bib de exemplu. [overwrite] e OBLIGATORIU exact din
%      motivul asta.
%
%   3. O singura adaugire de preambul dupa \usepackage{graphicx}: booktabs,
%      pentru \toprule / \midrule / \bottomrule din cele 4 tabele. NU se
%      adauga \bibliographystyle{acl_natbib}; acl.sty il seteaza singur si
%      duplicarea da eroare.
%
%   4. Lucrarea asta NU are figuri, deci nu are nici \graphicspath. graphicx
%      ramane fiindca e in template si e inofensiv.
%
%   5. Plasarea floats-urilor. Fiecare tabel e declarat inaintea titlului
%      care urmeaza primei lui referinte, nu in bloc la finalul sectiunii.
%      Declararea in bloc face LaTeX sa dedice coloane intregi numai lor.
%      Plus cele patru \renewcommand de fractii de float.
%
%   6. \author{} e gol. Anonimizarea o face template-ul: optiunea "review"
%      din acl.sty tipareste "Anonymous ACL submission" in locul lui \@author,
%      deci continutul lui nu se vede oricum la review. Gol e totusi mai bine
%      decit placeholderii, care ar pleca asa la tipar daca cineva comuta pe
%      "final". DE COMPLETAT la camera-ready.
%
%   7. \section*{Acknowledgments} e SCOS, obligatoriu la review conform
%      formatting.md. De repus la camera-ready, imediat inaintea referintelor.
%
% CE E DE VERIFICAT INAINTE DE TRIMITERE, si e scris si in locurile lui:
%   - lista de autori si subtitlul intrarii enomoto-etal-2024-tmu;
%   - citarile rapoartelor tehnice ale celor trei modele de limbaj, care
%     sunt acum numite in text fara citare;
%   - graful de citari MLSP 2024, pentru afirmatia din Related Work;
%   - titlul de mai jos, care e provizoriu.
%
% CIFRELE. Toate vin din artefacte pe disc si niciuna nu e recalculata la
% scriere. Sursele, o data pentru tot fisierul:
%   output/fusion_lolo_adoptat.csv     cifrele per limba ale sistemului
%   output/e_rezumat_modele.csv        jurnalul generat, cu diferente perechi
%   output/e_rezumat_modele_pre-p11_n5606.csv   arhiva, pentru Tabelul 4
%   output/e_adoptie.md                cele doua parghii, §1 si §2
%   output/e_m1.md                     alegerea mediaza-vs-selecteaza, LOLO
%   output/e_m3.md                     leave-one-encoder-out
%   output/e_pick_split.md             latimile criteriului, harta de modele
%   output/e_bakeoff_test.md           r per limba al fiecarui model, semnalul B
%   output/e_lemma.md                  fixul de lema si teza morfologica
%   output/repere_publicate.md         reperele publicate, transcrise de mana
%
% ATENTIE LA UN LUCRU CARE SE POATE SCRIE GRESIT FARA SA SE VADA: cifrele
% Tabelului 4 sunt pe n = 5606, dinaintea fixului de potrivire a spanului,
% iar restul lucrarii e pe n = 5624. Nu se scad unele din altele. Jurnalul
% marcheaza randurile vechi cu "n vechi (5606), NU se compara".
% =====================================================================

\begin{filecontents*}[overwrite]{custom.bib}
% Use this file for citations not found in the ACL Anthology (contained in "anthology.bib").

% Cele 7 intrari ale template-ului ACL. Raman necitate, dar inofensive:
% fisierul merge lipit si intr-un proiect care are deja un custom.bib.

@book{Aho:72,
    author  = {Alfred V. Aho and Jeffrey D. Ullman},
    title   = {The Theory of Parsing, Translation and Compiling},
    year    = "1972",
    volume  = "1",
    publisher = {Prentice-Hall},
    address = {Englewood Cliffs, NJ}
}

@book{APA:83,
    author  = {{American Psychological Association}},
    title   = {Publications Manual},
    year    = "1983",
    publisher = {American Psychological Association},
    address = {Washington, DC}
}

@article{Chandra:81,
	author = {Ashok K. Chandra and Dexter C. Kozen and Larry J. Stockmeyer},
	year = "1981",
	title = {Alternation},
	journal = {Journal of the Association for Computing Machinery},
	volume = "28",
	number = "1",
	pages = "114--133",
	doi = "10.1145/322234.322243",
}

@inproceedings{andrew2007scalable,
  title={Scalable training of {L1}-regularized log-linear models},
  author={Andrew, Galen and Gao, Jianfeng},
  booktitle={Proceedings of the 24th International Conference on Machine Learning},
  pages={33--40},
  year={2007},
}

@book{Gusfield:97,
    author  = {Dan Gusfield},
    title   = {Algorithms on Strings, Trees and Sequences},
    year    = "1997",
    publisher = {Cambridge University Press},
    address = {Cambridge, UK}
}

@article{rasooli-tetrault-2015,
    author    = {Mohammad Sadegh Rasooli and Joel R. Tetreault},
    title     = {Yara Parser: {A} Fast and Accurate Dependency Parser},
    journal   = {Computing Research Repository},
    volume    = {arXiv:1503.06733},
    year      = {2015},
    url       = {http://arxiv.org/abs/1503.06733},
    note    = {version 2}
}

@article{Ando2005,
	Acmid = {1194905},
	Author = {Ando, Rie Kubota and Zhang, Tong},
	Issn = {1532-4435},
	Issue_Date = {12/1/2005},
	Journal = {Journal of Machine Learning Research},
	Month = dec,
	Numpages = {37},
	Pages = {1817--1853},
	Publisher = {JMLR.org},
	Title = {A Framework for Learning Predictive Structures from Multiple Tasks and Unlabeled Data},
	Volume = {6},
	Year = {2005}
}

% ---- Intrarile copiate din Paper/Articol_Complexity/paper_full.tex, unde
% ---- au fost deja verificate contra ACL Anthology. Nu se re-transcriu.

@inproceedings{shardlow-etal-2021-semeval,
    title = "{S}em{E}val-2021 Task 1: Lexical Complexity Prediction",
    author = "Shardlow, Matthew  and
      Evans, Richard  and
      Paetzold, Gustavo Henrique  and
      Zampieri, Marcos",
    editor = "Palmer, Alexis  and
      Schneider, Nathan  and
      Schluter, Natalie  and
      Emerson, Guy  and
      Herbelot, Aurelie  and
      Zhu, Xiaodan",
    booktitle = "Proceedings of the 15th International Workshop on Semantic Evaluation (SemEval-2021)",
    month = aug,
    year = "2021",
    address = "Online",
    publisher = "Association for Computational Linguistics",
    url = "https://aclanthology.org/2021.semeval-1.1/",
    doi = "10.18653/v1/2021.semeval-1.1",
    pages = "1--16"
}

@inproceedings{yimam-etal-2018-report,
    title = "A Report on the Complex Word Identification Shared Task 2018",
    author = {Yimam, Seid Muhie  and
      Biemann, Chris  and
      Malmasi, Shervin  and
      Paetzold, Gustavo  and
      Specia, Lucia  and
      {\v{S}}tajner, Sanja  and
      Tack, Ana{\"i}s  and
      Zampieri, Marcos},
    editor = "Tetreault, Joel  and
      Burstein, Jill  and
      Kochmar, Ekaterina  and
      Leacock, Claudia  and
      Yannakoudakis, Helen",
    booktitle = "Proceedings of the Thirteenth Workshop on Innovative Use of {NLP} for Building Educational Applications",
    month = jun,
    year = "2018",
    address = "New Orleans, Louisiana",
    publisher = "Association for Computational Linguistics",
    url = "https://aclanthology.org/W18-0507/",
    doi = "10.18653/v1/W18-0507",
    pages = "66--78"
}

@inproceedings{shardlow-etal-2024-multilingual,
    title = "Multilingual Resources for Lexical Complexity Prediction: A Review",
    author = "Shardlow, Matthew  and
      North, Kai  and
      Zampieri, Marcos",
    editor = "Nunzio, Giorgio Maria Di  and
      Vezzani, Federica  and
      Ermakova, Liana  and
      Azarbonyad, Hosein  and
      Kamps, Jaap",
    booktitle = "Proceedings of the Workshop on DeTermIt! Evaluating Text Difficulty in a Multilingual Context @ LREC-COLING 2024",
    month = may,
    year = "2024",
    address = "Torino, Italia",
    publisher = "ELRA and ICCL",
    url = "https://aclanthology.org/2024.determit-1.5/",
    pages = "51--59"
}

@inproceedings{shardlow-etal-2024-bea,
    title = "The {BEA} 2024 Shared Task on the Multilingual Lexical Simplification Pipeline",
    author = {Shardlow, Matthew  and
      Alva-Manchego, Fernando  and
      Batista-Navarro, Riza  and
      Bott, Stefan  and
      Calderon Ramirez, Saul  and
      Cardon, R{\'e}mi  and
      Fran{\c{c}}ois, Thomas  and
      Hayakawa, Akio  and
      Horbach, Andrea  and
      H{\"u}lsing, Anna  and
      Ide, Yusuke  and
      Imperial, Joseph Marvin  and
      Nohejl, Adam  and
      North, Kai  and
      Occhipinti, Laura  and
      Rojas, Nelson Per{\'e}z  and
      Raihan, Nishat  and
      Ranasinghe, Tharindu  and
      Salazar, Martin Solis  and
      {\v{S}}tajner, Sanja  and
      Zampieri, Marcos  and
      Saggion, Horacio},
    editor = {Kochmar, Ekaterina  and
      Bexte, Marie  and
      Burstein, Jill  and
      Horbach, Andrea  and
      Laarmann-Quante, Ronja  and
      Tack, Ana{\"i}s  and
      Yaneva, Victoria  and
      Yuan, Zheng},
    booktitle = "Proceedings of the 19th Workshop on Innovative Use of NLP for Building Educational Applications (BEA 2024)",
    month = jun,
    year = "2024",
    address = "Mexico City, Mexico",
    publisher = "Association for Computational Linguistics",
    url = "https://aclanthology.org/2024.bea-1.51/",
    pages = "571--589"
}

@inproceedings{saggion-etal-2024-lexical,
    title = "Lexical Complexity Prediction and Lexical Simplification for {C}atalan and {S}panish: Resource Creation, Quality Assessment, and Ethical Considerations",
    author = "Saggion, Horacio  and
      Bott, Stefan  and
      Szasz, Sandra  and
      P{\'e}rez, Nelson  and
      Calder{\'o}n, Sa{\'u}l  and
      Sol{\'i}s, Mart{\'i}n",
    editor = "Shardlow, Matthew  and
      Saggion, Horacio  and
      Alva-Manchego, Fernando  and
      Zampieri, Marcos  and
      North, Kai  and
      {\v{S}}tajner, Sanja  and
      Stodden, Regina",
    booktitle = "Proceedings of the Third Workshop on Text Simplification, Accessibility and Readability (TSAR 2024)",
    month = nov,
    year = "2024",
    address = "Miami, Florida, USA",
    publisher = "Association for Computational Linguistics",
    url = "https://aclanthology.org/2024.tsar-1.9/",
    doi = "10.18653/v1/2024.tsar-1.9",
    pages = "82--94"
}

@inproceedings{qiu-etal-2024-complex,
    title = "{C}omp{L}ex-{ZH}: A New Dataset for Lexical Complexity Prediction in {M}andarin and {C}antonese",
    author = "Qiu, Le  and
      Guo, Shanyue  and
      Wong, Tak-Sum  and
      Chersoni, Emmanuele  and
      Lee, John  and
      Huang, Chu-Ren",
    editor = "Shardlow, Matthew  and
      Saggion, Horacio  and
      Alva-Manchego, Fernando  and
      Zampieri, Marcos  and
      North, Kai  and
      {\v{S}}tajner, Sanja  and
      Stodden, Regina",
    booktitle = "Proceedings of the Third Workshop on Text Simplification, Accessibility and Readability (TSAR 2024)",
    month = nov,
    year = "2024",
    address = "Miami, Florida, USA",
    publisher = "Association for Computational Linguistics",
    url = "https://aclanthology.org/2024.tsar-1.3/",
    doi = "10.18653/v1/2024.tsar-1.3",
    pages = "20--26"
}

@inproceedings{cristea-nisioi-2024-machine,
    title = "Archaeology at {MLSP} 2024: Machine Translation for Lexical Complexity Prediction and Lexical Simplification",
    author = "Cristea, Petru  and
      Nisioi, Sergiu",
    editor = {Kochmar, Ekaterina  and
      Bexte, Marie  and
      Burstein, Jill  and
      Horbach, Andrea  and
      Laarmann-Quante, Ronja  and
      Tack, Ana{\"i}s  and
      Yaneva, Victoria  and
      Yuan, Zheng},
    booktitle = "Proceedings of the 19th Workshop on Innovative Use of NLP for Building Educational Applications (BEA 2024)",
    month = jun,
    year = "2024",
    address = "Mexico City, Mexico",
    publisher = "Association for Computational Linguistics",
    url = "https://aclanthology.org/2024.bea-1.55/",
    pages = "610--617"
}

@inproceedings{conneau2020xlmr,
  title        = {Unsupervised Cross-lingual Representation Learning at Scale},
  author       = {Conneau, Alexis and Khandelwal, Kartikay and Goyal, Naman and Chaudhary, Vishrav and Wenzek, Guillaume and Guzm{\'a}n, Francisco and Grave, Edouard and Ott, Myle and Zettlemoyer, Luke and Stoyanov, Veselin},
  booktitle    = {Proceedings of the 58th Annual Meeting of the Association for Computational Linguistics (ACL)},
  pages        = {8440--8451},
  year         = {2020},
  url          = {https://aclanthology.org/2020.acl-main.747}
}

@inproceedings{pires-etal-2019-multilingual,
    title = "How Multilingual is Multilingual {BERT}?",
    author = "Pires, Telmo  and
      Schlinger, Eva  and
      Garrette, Dan",
    editor = "Korhonen, Anna  and
      Traum, David  and
      M{\`a}rquez, Llu{\'i}s",
    booktitle = "Proceedings of the 57th Annual Meeting of the Association for Computational Linguistics",
    month = jul,
    year = "2019",
    address = "Florence, Italy",
    publisher = "Association for Computational Linguistics",
    url = "https://aclanthology.org/P19-1493/",
    doi = "10.18653/v1/P19-1493",
    pages = "4996--5001"
}

@inproceedings{wu-dredze-2020-languages,
    title = "Are All Languages Created Equal in Multilingual {BERT}?",
    author = "Wu, Shijie  and
      Dredze, Mark",
    editor = "Gella, Spandana  and
      Welbl, Johannes  and
      Rei, Marek  and
      Petroni, Fabio  and
      Lewis, Patrick  and
      Strubell, Emma  and
      Seo, Minjoon  and
      Hajishirzi, Hannaneh",
    booktitle = "Proceedings of the 5th Workshop on Representation Learning for NLP",
    month = jul,
    year = "2020",
    address = "Online",
    publisher = "Association for Computational Linguistics",
    url = "https://aclanthology.org/2020.repl4nlp-1.16/",
    doi = "10.18653/v1/2020.repl4nlp-1.16",
    pages = "120--130"
}

@inproceedings{peters-etal-2019-tune,
    title = "To Tune or Not to Tune? Adapting Pretrained Representations to Diverse Tasks",
    author = "Peters, Matthew E.  and
      Ruder, Sebastian  and
      Smith, Noah A.",
    editor = "Augenstein, Isabelle  and
      Gella, Spandana  and
      Ruder, Sebastian  and
      Kann, Katharina  and
      Can, Burcu  and
      Welbl, Johannes  and
      Conneau, Alexis  and
      Ren, Xiang  and
      Rei, Marek",
    booktitle = "Proceedings of the 4th Workshop on Representation Learning for NLP (RepL4NLP-2019)",
    month = aug,
    year = "2019",
    address = "Florence, Italy",
    publisher = "Association for Computational Linguistics",
    url = "https://aclanthology.org/W19-4302/",
    doi = "10.18653/v1/W19-4302",
    pages = "7--14"
}

@inproceedings{gorman-bedrick-2019-need,
    title = "We Need to Talk about Standard Splits",
    author = "Gorman, Kyle  and
      Bedrick, Steven",
    editor = "Korhonen, Anna  and
      Traum, David  and
      M{\`a}rquez, Llu{\'i}s",
    booktitle = "Proceedings of the 57th Annual Meeting of the Association for Computational Linguistics",
    month = jul,
    year = "2019",
    address = "Florence, Italy",
    publisher = "Association for Computational Linguistics",
    url = "https://aclanthology.org/P19-1267/",
    doi = "10.18653/v1/P19-1267",
    pages = "2786--2791"
}

% ---- Intrarile NOI, scrise pentru lucrarea multilingva. Regulile din
% ---- memoria `paper-overleaf-workflow`: zero Unicode, zero arond in
% ---- comentarii, acronimele in acolade ca .bst-ul sa nu le scrie cu
% ---- minuscule.

% Sursa reperelor publicate. Datele bibliografice sunt cele verificate in
% output/repere_publicate.md, unde e transcris si Tabelul 13.
@article{nohejl-etal-2025-japanese,
    author  = {Nohejl, Adam and Hayakawa, Akio and Ide, Yusuke and Watanabe, Taro},
    title   = {A {J}apanese Dataset and Efficient Multilingual {LLM}-Based Methods for Lexical Simplification and Lexical Complexity Prediction},
    journal = {Journal of Natural Language Processing},
    volume  = {32},
    number  = {4},
    pages   = {1129--1163},
    year    = {2025}
}

% DE COMPLETAT inainte de trimitere: lista completa de autori si subtitlul,
% de pe aclanthology.org/2024.bea-1.52.bib. Ce e mai jos e numai ce e
% verificat in output/repere_publicate.md, si "and others" randeaza corect
% "Enomoto et al." pana atunci.
@inproceedings{enomoto-etal-2024-tmu,
    title     = "{TMU}-{HIT} at {MLSP} 2024",
    author    = "Enomoto and others",
    booktitle = "Proceedings of the 19th Workshop on Innovative Use of NLP for Building Educational Applications (BEA 2024)",
    month     = jun,
    year      = "2024",
    address   = "Mexico City, Mexico",
    publisher = "Association for Computational Linguistics",
    url       = "https://aclanthology.org/2024.bea-1.52/"
}

@inproceedings{liu-etal-2023-geval,
    title     = "{G}-{E}val: {NLG} Evaluation using {GPT}-4 with Better Human Alignment",
    author    = "Liu, Yang  and
      Iter, Dan  and
      Xu, Yichong  and
      Wang, Shuohang  and
      Xu, Ruochen  and
      Zhu, Chenguang",
    booktitle = "Proceedings of the 2023 Conference on Empirical Methods in Natural Language Processing",
    month     = dec,
    year      = "2023",
    address   = "Singapore",
    publisher = "Association for Computational Linguistics",
    url       = "https://aclanthology.org/2023.emnlp-main.153/",
    doi       = "10.18653/v1/2023.emnlp-main.153",
    pages     = "2511--2522"
}

@inproceedings{feng-etal-2022-language,
    title     = "Language-agnostic {BERT} Sentence Embedding",
    author    = "Feng, Fangxiaoyu  and
      Yang, Yinfei  and
      Cer, Daniel  and
      Arivazhagan, Naveen  and
      Wang, Wei",
    booktitle = "Proceedings of the 60th Annual Meeting of the Association for Computational Linguistics (Volume 1: Long Papers)",
    month     = may,
    year      = "2022",
    address   = "Dublin, Ireland",
    publisher = "Association for Computational Linguistics",
    url       = "https://aclanthology.org/2022.acl-long.62/",
    doi       = "10.18653/v1/2022.acl-long.62",
    pages     = "878--891"
}

@article{wang2024multilinguale5,
    author  = {Wang, Liang and Yang, Nan and Huang, Xiaolong and Yang, Linjun and Majumder, Rangan and Wei, Furu},
    title   = {Multilingual {E}5 Text Embeddings: A Technical Report},
    journal = {Computing Research Repository},
    volume  = {arXiv:2402.05672},
    year    = {2024},
    url     = {https://arxiv.org/abs/2402.05672}
}

@inproceedings{qi-etal-2020-stanza,
    title     = "{S}tanza: A {P}ython Natural Language Processing Toolkit for Many Human Languages",
    author    = "Qi, Peng  and
      Zhang, Yuhao  and
      Zhang, Yuhui  and
      Bolton, Jason  and
      Manning, Christopher D.",
    booktitle = "Proceedings of the 58th Annual Meeting of the Association for Computational Linguistics: System Demonstrations",
    month     = jul,
    year      = "2020",
    address   = "Online",
    publisher = "Association for Computational Linguistics",
    url       = "https://aclanthology.org/2020.acl-demos.14/",
    doi       = "10.18653/v1/2020.acl-demos.14",
    pages     = "101--108"
}

@inproceedings{rust-etal-2021-good,
    title     = "How Good is Your Tokenizer? On the Monolingual Performance of Multilingual Language Models",
    author    = "Rust, Phillip  and
      Pfeiffer, Jonas  and
      Vuli{\'c}, Ivan  and
      Ruder, Sebastian  and
      Gurevych, Iryna",
    booktitle = "Proceedings of the 59th Annual Meeting of the Association for Computational Linguistics and the 11th International Joint Conference on Natural Language Processing (Volume 1: Long Papers)",
    month     = aug,
    year      = "2021",
    address   = "Online",
    publisher = "Association for Computational Linguistics",
    url       = "https://aclanthology.org/2021.acl-long.243/",
    doi       = "10.18653/v1/2021.acl-long.243",
    pages     = "3118--3135"
}
\end{filecontents*}

\documentclass[11pt]{article}

% Change "review" to "final" to generate the final (sometimes called camera-ready) version.
% Change to "preprint" to generate a non-anonymous version with page numbers.
\usepackage[review]{acl}

% Standard package includes
\usepackage{times}
\usepackage{latexsym}

% For proper rendering and hyphenation of words containing Latin characters (including in bib files)
\usepackage[T1]{fontenc}
% For Vietnamese characters
% \usepackage[T5]{fontenc}
% See https://www.latex-project.org/help/documentation/encguide.pdf for other character sets

% This assumes your files are encoded as UTF8
\usepackage[utf8]{inputenc}

% This is not strictly necessary, and may be commented out,
% but it will improve the layout of the manuscript,
% and will typically save some space.
\usepackage{microtype}

% This is also not strictly necessary, and may be commented out.
% However, it will improve the aesthetics of text in
% the typewriter font.
\usepackage{inconsolata}

%Including images in your LaTeX document requires adding
%additional package(s)
\usepackage{graphicx}

% ---------------------------------------------------------------------
% Abaterea 3: singura adaugire de preambul ceruta de corpul lucrarii.
% ---------------------------------------------------------------------
\usepackage{booktabs}          % \toprule / \midrule / \bottomrule in the four tables

% Fractiile implicite ale lui LaTeX (topfraction 0.7, textfraction 0.2)
% dedica coloane intregi numai tabelelor. Relaxate ca sa incapa text in
% aceeasi coloana cu un float.
\renewcommand{\topfraction}{0.9}
\renewcommand{\bottomfraction}{0.5}
\renewcommand{\textfraction}{0.08}
\renewcommand{\floatpagefraction}{0.8}

% If the title and author information does not fit in the area allocated, uncomment the following
%
%\setlength\titlebox{<dim>}
%
% and set <dim> to something 5cm or larger.

% TITLUL E PROVIZORIU. Numeste ambele contributii, cea de protocol si cea de
% rezultat. De confirmat cu coordonatorul inainte de trimitere.
\title{Lexical Complexity Prediction Without Target-Language Labels:\\
       When to Select a Component and When to Average}

% ---------------------------------------------------------------------
% Blocul de autori e GOL, intentionat. Vezi abaterea 6 din capul fisierului.
%
% DE COMPLETAT la camera-ready, dupa modelul template-ului:
%   \author{Nume Autor \\
%     Afiliere \\
%     \texttt{email@domain} \\\And
%     Al Doilea Autor \\
%     Afiliere \\
%     \texttt{email@domain} \\}
% ---------------------------------------------------------------------
\author{}

\begin{document}
\maketitle

\begin{abstract}
We predict lexical complexity in ten languages without using a single
annotated example from the language being evaluated. Two signals are
fused: a regressor over a frozen multilingual encoder, fitted on
annotations from the other languages, and a complexity rating read from
the token probabilities of an open-weight language model run locally.
Every choice the pipeline makes, including which language model to use
and how to weight the two signals against each other, is fitted on the
other nine languages. Averaged over the ten, we reach a Pearson
correlation of $0.675$, ahead by $0.065$ of the strongest published
result obtained without target-language labels, and ahead on nine of the
ten languages. We further show that the same leave-one-language-out
criterion can decide whether a component should be selected per language
or averaged over all candidates, and that it answers that question in
opposite ways for the two components of one pipeline.
\end{abstract}

% =====================================================================
% Sectiunea 1, Introduction. Scrisa ULTIMA, cum cere ordinea.
%
% Structura ceruta de coordonator, ghidul Columbia hgs, vezi memoria
% `intro-style-hgs`. Maparea celor 5 puncte pe cele 6 paragrafe:
%   P1 motivatie          -> ¶1
%   P2 problema specifica -> ¶2 + ¶3
%   P3 contributiile      -> ¶4, paragraf narativ, persoana I, fara buline
%   P4 diferentierea      -> ¶5, la nivel inalt, se inchide cu trimitere la §2
%   P5 roadmap            -> ¶6, enumera CONTINUT, nu numere de sectiune
%
% Reguli aplicate: fara deschidere de tip definitie de dictionar, fara
% aforisme, fara formule, fara cifre de sprijin peste cifra de titlu, fraza
% -tema la inceputul fiecarui paragraf, timpuri verbale consistente in
% interiorul paragrafului.
%
% Contributia a doua e REFORMULATA fata de pasajul 10 din
% Paper/TEXT_DE_LIPIT_multilingv.tex: acolo scria „recovers most of the gap",
% ceea ce supraliciteaza. Decalajul francez fata de cea mai mare cifra
% publicata a trecut de la 0.192 la 0.129, adica o treime.
% =====================================================================
\section{Introduction}
\label{sec:intro}

Deciding which words in a text will be difficult is the first step of
automatic simplification, and it bounds everything the rest of the
pipeline is able to do, since a word that is never identified as hard is
never replaced. The readers this matters for are not one group.
Language learners, readers with limited literacy, and specialists
reading outside their own field each meet the problem in a different
vocabulary, and a system that serves one of them does not thereby serve
the others. Lexical complexity prediction estimates how difficult a
given word is in the sentence where it occurs, which is what lets a
simplification system rank candidates for replacement and a readability
measure locate where in a text the difficulty sits.

Nearly every reported system for the task is built the same way. Human
annotators rate words in context in some language, an encoder is
fine-tuned on those ratings, and the resulting model predicts in that
language. This works, and it has produced the figures the field cites,
but it makes the method available exactly where it is least needed. The
languages with complexity annotations are the languages that already
have simplification tools, readability formulas and graded reading
material; the languages without them have none of these, and no obvious
route to any.

Removing the labels breaks more than the training step. It also removes
the means by which every other decision is normally made. Which
pretrained encoder suits a language, which of several language models
rates its words most reliably, how far to trust a transfer model against
a direct one: all of these are ordinarily settled by validating in the
target language, and none of them can be settled that way when the
target language is the one thing that must not be looked at. The surface
features carry a second, quieter version of the problem, since character
length and syllable count do not mean the same thing across writing
systems, and frequency lists index inflected forms as though each were
its own word.

We take the position that these decisions should be made by the same
procedure that produces the predictions, and we report what happens when
they are. First, we show that the choice between selecting a component
per language and averaging over all candidates can be made by the same
leave-one-language-out criterion that fixes every other choice, and that
the criterion gives opposite answers for two components of one pipeline,
each answer confirmed by measurement. Second, we show that fitting one
component on the trial split while fitting the rest on the evaluation
split is a measurable error and not a detail of bookkeeping, and that
correcting it closes a third of the gap on the language where the gap to
published results was largest. Third, we replace character length with
tokeniser fertility and resolve frequency lookups through lemmas, two
changes argued from morphology rather than chosen by score, which
transfer across writing systems. The resulting predictions reach a
Pearson correlation of $0.675$ averaged over ten languages, using no annotated
example from the language being evaluated, which is ahead by $0.065$ of
the strongest published result obtained under that same condition.
Fourth, we report every configuration we measured, together with what
choosing between them cost us.

Two things separate this from the closest published work. The strongest
published averages on this data come from prompting a proprietary model
with labelled examples in the target language, so they answer a
different question than the one we ask, and the distance between those
figures and ours is smaller in either direction than the width of our
own confidence interval.
The zero-label results that do answer our question were obtained by
translating into English, which concedes the multilingual problem
instead of addressing it. Section~\ref{sec:related} places both, along
with the transfer and tokenisation work the approach builds on.

The rest of the paper describes the two signals and the criterion that
combines them, reports per-language correlations against every published
figure on this test set, separates what each of the two corrections
contributes, and states where the reported configuration was chosen by
us rather than by the procedure.

% =====================================================================
% Sectiunea 2, Related Work. Pe \paragraph{}, fara subsectiuni numerotate,
% ca in RALS. Impreuna cu introducerea trebuie sa stea sub 2 pagini.
%
% Fiecare bloc se inchide cu o fraza de gap, tiparul luat din
% 2605.07466_6048.pdf. Nicio citare nu apare in doua blocuri: regula
% „nicio idee spusa de doua ori" prinde exact repetitia de citari.
%
% NEVERIFICAT, si trebuie verificat inainte de trimitere: graful de citari
% al datelor MLSP 2024, pentru afirmatia ca nimeni n-a raportat protocolul
% asta pe setul asta. Formularea de mai jos spune ce am gasit noi, nu ca nu
% exista, tocmai fiindca verificarea nu s-a facut.
% =====================================================================
\section{Related Work}
\label{sec:related}

\paragraph{Lexical complexity prediction.} Predicting how difficult a
word is for a reader, in the context where it occurs, was posed as a
regression problem by \citet{shardlow-etal-2021-semeval} for English,
after earlier work treated it as binary identification
\citep{yimam-etal-2018-report}. Data has since appeared for other
languages one at a time, for Catalan and Spanish
\citep{saggion-etal-2024-lexical}, for Mandarin and Cantonese
\citep{qiu-etal-2024-complex}, and, in the collection we use here, for
ten languages annotated under a shared protocol
\citep{shardlow-etal-2024-bea}. \citet{shardlow-etal-2024-multilingual}
survey what exists. The resources are what make a multilingual
evaluation possible, and they are also what makes the labelled-data
assumption easy to keep: with annotations available in ten languages, a
system can be fitted in each of them and never asked what it would do in
the eleventh.

\paragraph{Systems on this data.} Reported approaches divide by what they
condition on. \citet{cristea-nisioi-2024-machine} translate into English
and apply an English model, using no target-language labels and reaching
$0.486$. TMU-HIT \citep{enomoto-etal-2024-tmu} prompt a large
proprietary model with labelled examples in the target language, and
their zero-shot run is the strongest published result that uses none.
\citet{nohejl-etal-2025-japanese} give the current best figures, by
reading complexity from the token probabilities of a proprietary model
rather than from its generated text, following the evaluation
formulation of \citet{liu-etal-2023-geval}, and adding a temperature
fitted on the trial sets of all ten languages. What is missing from this
picture is a system that reaches that range while conditioning on
nothing in the target language and running on weights that can be
downloaded.

\paragraph{Transfer through frozen encoders.} Multilingual encoders
transfer across languages without parallel supervision
\citep{pires-etal-2019-multilingual, conneau2020xlmr}, unevenly, and
worst for the languages with least pretraining data
\citep{wu-dredze-2020-languages}. \citet{peters-etal-2019-tune} show
that freezing rather than fine-tuning is often the better choice when
the target task has little supervision, which is the regime here, since
by construction there is none. What is less settled is how to choose
among encoders when no labels exist in the language to choose for, and
this is the decision we take from the researcher and give to the
protocol.

\paragraph{Frequency, morphology and tokenisation.} Word frequency is
the oldest predictor of lexical difficulty and remains the strongest
single feature, but it is measured over surface forms, which in an
inflecting language splits one word's evidence across many entries. The
number of subword units a tokeniser assigns to a word varies sharply
across languages for the same encoder \citep{rust-etal-2021-good}, which
makes it both a symptom of uneven coverage and, as we use it here, a
feature that is measured in the same unit in every language. Neither
observation appears to have been carried into complexity prediction,
where character length is still the standard surface feature.

% =====================================================================
% Sectiunea 3, Approach. Scrisa DUPA Results, cum cere ordinea.
%
% Fiecare detaliu de aici a fost verificat in cod in sesiunea din 19 august:
%   trasaturile lexicale        solution_bert.py, `lexical`
%   fertilitatea                solution_bert.py, `fertility`
%   centrarea si concatenarea   solution_bert.py, `_Xc`
%   regresorul                  solution_bert.py, `blend_ridge_hgb`
%   cele doua prompturi         solution_bert.py:383-390
%   citirea cifrelor 1..5       solution_bert.py, `digit_mass` si `_score_msgs`
%   nivelurile de selectie      output/e_fusion_lolo_adoptat.md, tabelul de audit
%   latimile criteriilor        output/e_pick_split.md, output/e_encoder.md
%   M1                          output/e_m1.md
%   scara bratului A            output/e_rezumat_modele_pre-p11_n5606.csv
%
% Modelele de limbaj sunt numite in text FARA citare: rapoartele lor tehnice
% nu au fost verificate in sesiunea asta si nu se inventeaza intrari de
% bibliografie. DE COMPLETAT inainte de trimitere.
% =====================================================================
\section{Approach}
\label{sec:approach}

\subsection{Data}
\label{sec:data}

We use the multilingual lexical complexity data of
\citet{shardlow-etal-2024-bea}, which covers Catalan, English, Filipino,
German, French, Italian, Japanese, Portuguese, Sinhala and Spanish, with
a complexity value in $[0,1]$ for a target word in a sentence. We
evaluate on the test split of all ten languages, 5{,}624 instances.
English complexity annotations from CompLex
\citep{shardlow-etal-2021-semeval} are used as an additional labelled
source.

Throughout, the language being evaluated contributes no labels to
anything: not to training, not to feature selection, not to the choice
of language model, and not to the fusion weight. Section~\ref{sec:lolo}
states where each of those exclusions happens.

\subsection{The transfer signal}
\label{sec:signalA}

The first signal is a regressor over a frozen multilingual encoder. For
a target word we take the per-token hidden states of the last four
layers, restricted to the target span, and centre them on the mean of
the language. To this we concatenate four lexical columns, Zipf
frequency, a length measure, syllable count and capitalisation, plus two
indicators that record whether frequency and Latin-script surface
features are defined for that language at all. Sinhala has no frequency
list and Japanese has no meaningful syllable or capitalisation feature,
and marking those cases is what keeps a missing value from being read as
a small one. The regressor is the mean of a ridge regression and a
gradient-boosted tree, which is more stable across languages than either
alone.

Two of the lexical columns are not the obvious ones, and both changes
are motivated by morphology rather than chosen by score.

\paragraph{Frequency through the lemma.} An inflected form is often
rarer than the word it belongs to, and in a morphologically rich
language the frequency list will report that rarity as difficulty. We
therefore resolve frequency as the maximum of the Zipf frequency of the
surface form and that of its lemma, with lemmas from
\citet{qi-etal-2020-stanza}. The correction can only apply to the eight
languages that have both a lemmatiser and a frequency list. Across those
eight, what it gains correlates at $+0.78$ with a corpus measure of
morphological synthesis, which is the direction the argument predicts.
With eight points we report the coefficient and leave it at that.

\paragraph{Fertility instead of character length.} Character length is
not comparable across writing systems. A Japanese target in this data
averages $3.59$ characters and an Italian one $8.78$, so the same
numeric value means different things in the two languages. The number of
subword units the tokeniser produces for the target is by construction
in the same unit everywhere, and it is available without any model
weights or inference \citep{rust-etal-2021-good}. We use it in place of
character length. On Sinhala, where no frequency value exists, it is the
only frequency-like information that remains.

Table~\ref{tab:aarm} shows what the two changes contribute on their own.

% ---------------------------------------------------------------------
% Tabelul `tab:aarm`. Scara bratului A, din arhiva
% output/e_rezumat_modele_pre-p11_n5606.csv, randurile de brat inghetat:
%   0.5523 xlmr inghetat
%   0.5630 + lema
%   0.5753 + fertilitate in loc de lungime
%   0.5804 + ansamblu de surse
% TOATE pe n = 5606, adica DINAINTE de fixul de potrivire a spanului. Nu se
% scad din cifrele Tabelului 1, care sunt pe 5624. Comparatiile din interiorul
% tabelului sunt insa valide: ambele brate ale fiecarei perechi sunt pe
% aceleasi randuri, cu bootstrap pereche.
% Cifra +0.0107 si corelatia MSP sunt din output/e_lemma.md, tot pe n vechi.
% Propunerea de re-rulare pe 5624 e in raportul de final al sesiunii.
% ---------------------------------------------------------------------
\begin{table}[t]
\centering
\small
\begin{tabular}{lc}
\toprule
Transfer signal alone & avg $r$ \\
\midrule
Frozen encoder, surface frequency      & .5523 \\
\quad + frequency through the lemma    & .5630 \\
\quad + fertility for character length & .5753 \\
\quad + second transfer arm            & .5804 \\
\bottomrule
\end{tabular}
\caption{The transfer signal on its own, without the language model.
These four figures were measured on 5{,}606 instances, before the
span-matching correction that brought the count to 5{,}624, so they are
comparable with each other but not with Table~\ref{tab:results}. Each
step is a paired comparison on identical rows.}
\label{tab:aarm}
\end{table}

\paragraph{Two transfer arms.} We fit the regressor twice. One arm sees
only the English annotations and predicts every language from them. The
other adds the rows of the other nine languages, refitted once per
held-out language. The two are standardised within the language and
mixed with a single weight chosen on a nine-point grid, so the pipeline
can decide per language how much to lean on a source that never saw this
data at all.

\paragraph{Encoders.} We ran four frozen encoders through this
arm: XLM-R base and XLM-R large \citep{conneau2020xlmr}, LaBSE
\citep{feng-etal-2022-language} and multilingual E5
\citep{wang2024multilinguale5}. No single one of them improves the
fused result over XLM-R base with an interval that excludes zero. What
works is averaging their standardised predictions, and
Section~\ref{sec:selectavg} explains why that is the outcome of the
protocol rather than a preference.

\subsection{The language model signal}
\label{sec:signalB}

The second signal is a complexity rating from an open-weight language
model run locally. We use two prompts, both in English and both naming
the target language explicitly, which ask for a rating from 1 to 5 for
the target word in its sentence. Rather than sampling a reply, we read
the probability mass the model assigns to the tokens \texttt{1} to
\texttt{5} at the next position, take the expectation under that
distribution, rescale it to $[0,1]$, and average over the two prompts.
This costs one forward pass per prompt and no generation. Reading
complexity from token probabilities in this way is not our contribution;
it is the method of \citet{nohejl-etal-2025-japanese}, and we adopt it.

Three models were run: Gemma~4 12B at 4-bit precision, Qwen3 8B at
4-bit, and Qwen3.5 9B at 8-bit. Which of them serves a given language is
decided by the same criterion as everything else, and on the evaluation
split it selects the same model for all ten languages.

\subsection{Fusion, and where each choice is made}
\label{sec:lolo}

Both signals are standardised within the language and combined with a
single weight on a twenty-one point grid. For a held-out language $L$,
every quantity that the prediction for $L$ depends on is fitted on the
other nine languages: the transfer regressor, the mixing weight between
the two transfer arms, the choice of language model, and the fusion
weight. The criterion that ranks candidates for $L$ is itself nested. A
candidate is scored by its mean correlation over the nine fitting
languages, and the fusion weight used at each of those nine comes from
the remaining eight, so no candidate is ever judged with a weight fitted
on the very language that measures it.

One asymmetry in this scheme was a genuine error, and correcting it is
the second of the two changes in Table~\ref{tab:ablation}. The fusion
weight was fitted on the evaluation split while the language model was
chosen from trial-split scores. Nothing motivated the split; the reason
that had once justified it for the weight had been removed and left
standing for the model. Under the corrected procedure both read the same
split, always excluding the held-out language. The effect is confined
almost entirely to French, which gains $0.0583$, because that is the
only language where the two splits disagreed about which model to use.
This is worth stating precisely: the claim is not that different
languages need different models, since on the correct split the same
model wins everywhere. The claim is that trial scores are a poor proxy
for test scores on this signal, and this is the second measurement that
shows it.

\subsection{Selecting or averaging}
\label{sec:selectavg}

The same kind of decision arises twice, once for the encoder and once
for the language model, and the two do not have the same answer. What
decides is how far apart the candidates are under the criterion that
ranks them, measured against how far apart the fusion weight curve
spreads over its own grid. That curve is the reference for what a
difference worth acting on looks like here, since it is the one
quantity we already fit by search and know to matter. Where the
candidates spread flat against that reference, the maximum is noise and
averaging is safer. Where they spread wide, the maximum means something
and selection is safer.

The two cases fall on opposite sides, and both verdicts are confirmed by
measurement. For the language model the criterion spans $0.2677$ against
a weight curve spanning $0.0888$, three times wider, and selection is
right: it improves French by $0.0583$. For the encoder the criterion
spans $0.0077$ against $0.082$, more than ten times flatter, and
selection is wrong: choosing an encoder per language under the same
protocol scores $0.0050$ below using one encoder everywhere, and it
picks the best encoder for a language zero times out of ten. Averaging the encoders instead lifts the transfer signal to the score
that picking the best encoder per language would have reached, and it
does so without using any labels.

That leaves one binary decision that we could have made after looking at
test scores, and we removed it from our hands. For each target language
we compare averaging against selection using only the other nine
languages, through the same criterion that already fixes the fusion
weight and the language model. Averaging is chosen on all ten folds,
with a minimum margin of $0.0088$.

The unanimity is not a surprise, and it is worth saying why. Averaging
beats selection on nine of the ten languages, so withholding one
language cannot reverse the maximum. What this establishes is not that a
close decision came out correctly, but that the decision is stable and
rests on no single language. A criterion that flipped when one language
was withheld would have been evidence that the choice was noise, which
is what we observe for encoder selection itself.

One consequence should be stated plainly. On Filipino, averaging is
worse than selection, and the procedure selects averaging there anyway,
because it is not allowed to look at Filipino. That is a real loss we
pay to keep the procedure clean.

Choosing which three of the four encoders to average is a separate
question, and the criterion above does not answer it, because the four
candidates are not distinguishable. We take it as an engineering
decision, on cost and on behaviour, and Section~\ref{sec:notmax} reports
what it costs in score.

% =====================================================================
% Sectiunea 4, Results. Scrisa PRIMA in sesiunea din 19 august, fiindca
% e complet determinata de masuratori.
%
% Fiecare cifra si de unde vine:
%   0.6749, per limba          output/fusion_lolo_adoptat.csv, coloana r_fuz_raw
%   0.664581 / 0.669223 /
%   0.669735 / 0.674946        output/e_rezumat_modele.csv + output/e_adoptie.md §1
%   0.0052 / 0.0057 / 0.0104   output/e_adoptie.md §1, tabelul comparatiilor perechi
%   106%                       output/e_adoptie.md §2
%   ponderile 0.390..0.430     output/e_adoptie.md §2, paragraful despre franceza
%   reperele publicate         output/repere_publicate.md
%   0.675996, +0.0011          output/e_rezumat_modele.csv, randul „r MAXIM masurat"
%   leave-one-encoder-out      output/e_m3.md
%
% ATENTIE, greseala prinsa in TEXT_DE_LIPIT_multilingv.tex pasajul 2: cele
% doua ablatii erau INVERSATE acolo. Scoaterea mediei de encodere costa
% 0.0052 (adoptat minus picktest singur), scoaterea selectiei pe test costa
% 0.0057 (adoptat minus encmean singur). Asa e in e_adoptie.md §1.
% =====================================================================

% ---------------------------------------------------------------------
% Tabelul `tab:results`. Cifrele noastre per limba vin din output/fusion_lolo_adoptat.csv,
% coloana r_fuz_raw, rotunjite la 3 zecimale:
%   ca .541096  en .870245  fil .512088  de .733289  fr .645033
%   it .664653  ja .783530  pt .842722  si .396198  es .760611
%   medie exacta 0.6749464 -> .675
% Reperele sunt din output/repere_publicate.md, transcrise de mana din
% Nohejl et al. 2025 Tabelul 13, si din Anexa B a lui Shardlow et al. 2024a
% pentru randul TMU-HIT zero-shot uniform, care NU exista publicat ca medie.
%
% DOUA corecturi fata de Paper/TEXT_DE_LIPIT_multilingv.tex pasajul 9:
%   1. Archaeology e .486, nu .487. Cifra lui Nohejl e 0.487, dar media
%      valorilor per limba TIPARITE in randul asta e 0.4864. Un reviewer
%      care aduna randul trebuie sa obtina ce scrie in coloana medie. Din
%      acelasi motiv scrie .486 si in Related Work, unde cifra e citata a
%      doua oara: doua valori diferite pentru acelasi sistem in aceeasi
%      lucrare se vad imediat.
%   2. Captionul de acolo se contrazicea: spunea „only ... use no labels"
%      si apoi „TMU-HIT zero-shot uses none either". Rescris intr-o singura
%      propozitie, cu numarul de randuri.
% ---------------------------------------------------------------------
\begin{table*}[t]
\centering
\small
\begin{tabular}{lccccccccccc}
\toprule
& ca & en & fil & de & fr & it & ja & pt & si & es & avg \\
\midrule
Frequency baseline         & .301 & .748 & .389 & .591 & .517 & .519 & .642 & .713 & -.196 & .551 & .478 \\
Archaeology                & .274 & .790 & .443 & .551 & .533 & .534 & .485 & .683 & .044 & .527 & .486 \\
TMU-HIT, zero-shot         & .528 & .804 & \textbf{.569} & .658 & .625 & .539 & .645 & .764 & .248 & .720 & .610 \\
\midrule
TMU-HIT, per-language best & \textbf{.616} & .820 & \textbf{.569} & .718 & .625 & .601 & .733 & .786 & .308 & .762 & .654 \\
G-Eval, GPT-4o             & .573 & .827 & .494 & .711 & .764 & .629 & .745 & .761 & .328 & .739 & .657 \\
G-Scale, GPT-4o            & .583 & .840 & .504 & .714 & \textbf{.774} & .654 & .767 & .774 & .330 & .749 & .669 \\
\midrule
Ours                       & .541 & \textbf{.870} & .512 & \textbf{.733} & .645 & \textbf{.665} & \textbf{.784} & \textbf{.843} & \textbf{.396} & \textbf{.761} & \textbf{.675} \\
\bottomrule
\end{tabular}
\caption{Pearson correlation per language on the MLSP test set, and the
unweighted mean over the ten languages. Bold marks the highest figure in
each column. The first three rows and ours use no labels in the language
being evaluated, while the middle three condition on three to thirty
labelled examples in it. Published figures are from
\citet{nohejl-etal-2025-japanese}, except the TMU-HIT zero-shot row,
which we obtained by averaging the per-language figures of that run in
Appendix B of \citet{shardlow-etal-2024-bea}, since no mean over
languages is published for it.}
\label{tab:results}
\end{table*}

\section{Results}
\label{sec:results}

Table~\ref{tab:results} reports Pearson correlation per language,
averaged without weighting across the ten languages, on 5{,}624 test
instances. The configuration described in Section~\ref{sec:approach}
reaches $0.6749$, with a bootstrap interval of $[0.660, 0.689]$ for the
mean.

Every interval we report resamples contexts rather than instances,
2{,}000 times. Seven of the ten languages contain about three target
words per context, so an instance-level bootstrap would be narrowest
exactly where the dependence between items is strongest. Differences
between configurations are paired: the same resample selects the same
contexts on both sides, so the interval belongs to the difference and
not to the sum of two uncertainties.

\subsection{What each component contributes}
\label{sec:ablation}

Two changes separate this configuration from the one we started with,
which used a single encoder and chose the language model on the trial
split. Table~\ref{tab:ablation} measures them separately and together.
Removing the encoder average costs $0.0052$, removing the
evaluation-split model choice costs $0.0057$, and removing both costs
$0.0104$, with a paired interval of $[+0.006, +0.015]$.

The two changes almost do not overlap. Their separate gains sum to
$0.0098$, and measured together they give $0.0104$, so nothing is lost
by applying both. Their per-language profiles explain why: averaging
encoders helps most on Sinhala ($+0.0176$) and Catalan ($+0.0130$),
while the model choice acts almost entirely on French ($+0.0583$) and
leaves three languages identical row by row. The one place where they
could have cancelled is the fusion weight, which reallocates when both
arms improve. It does move, from $0.390$ averaged over the ten
languages to $0.430$ when the frozen arm improves and to $0.375$ when
the language model arm improves, and it settles at $0.410$ when both
apply, but not far enough to absorb either gain.

Two languages fall. Filipino loses $0.0051$ and Japanese loses
$0.0051$, both from the encoder average and both with intervals that
include zero.

% ---------------------------------------------------------------------
% Tabelul `tab:ablation`. Cele patru configuratii, din output/e_rezumat_modele.csv
% (coloanele r_mediu, d_vs_adoptat, ci_lo, ci_hi) si output/e_adoptie.md §1.
%   0.664581  ens              d -0.010365  [-0.0146, -0.0060]
%   0.669223  encmean3         d -0.005724  [-0.0087, -0.0029]
%   0.669735  picktest         d -0.005211  [-0.0084, -0.0017]
%   0.674946  ADOPTAT
% Celula de referinta e LASATA GOALA, nu umpluta cu o linie: regula fara
% em-dash se aplica si in tabele.
% ---------------------------------------------------------------------
\begin{table}[t]
\centering
\small
\begin{tabular}{lcc}
\toprule
 & avg $r$ & vs.\ reported \\
\midrule
Single encoder, model                 & .6646 & $-.0104$ \\
\quad chosen on trial                 &       & $[-.0146, -.0060]$ \\
\addlinespace[2pt]
Encoder average only                  & .6692 & $-.0057$ \\
                                      &       & $[-.0087, -.0029]$ \\
\addlinespace[2pt]
Evaluation-split model                & .6697 & $-.0052$ \\
\quad choice only                     &       & $[-.0084, -.0017]$ \\
\addlinespace[2pt]
\textbf{Both, reported}               & \textbf{.6749} & \\
\bottomrule
\end{tabular}
\caption{The two changes, separately and together, on the same 5{,}624
test instances. Differences are paired against the reported
configuration, with 95\% bootstrap intervals over resampled contexts.}
\label{tab:ablation}
\end{table}

\subsection{Comparison with published figures}
\label{sec:published}

Two published systems set the frame. TMU-HIT run under a uniform
zero-shot protocol reaches $0.610$, and is the strongest published
result on this data that uses no labels in the language it is
evaluated on. We are ahead of it by $0.065$, winning on nine of the ten
languages. G-Scale with GPT-4o reaches $0.669$ and is the highest
published average on this test set. We are ahead of it by $0.006$.

That second margin should not be read as an improvement over the state
of the art, for three reasons. The interval of our own mean is
$[0.660, 0.689]$, so a difference of $0.006$ is a fifth of one interval
width. It cannot be tested, because per-instance predictions for
G-Scale are not available. And the protocols differ: G-Scale conditions
on thirty labelled examples in the target language and on a temperature
fitted across all ten trial sets, while we condition on none. Being
above it is therefore not a protocol win, just as being below it would
not have been a protocol failure. We keep it in the table because it is
the highest published figure and its absence would be noticed.

Under the comparison that does hold the protocol fixed, the margin is
$0.065$ rather than $0.006$, and seven of the ten per-language figures
we report are the highest published for that language. The three that
are not are Catalan, Filipino and French.

\subsection{The figure we report is not the highest we measured}
\label{sec:notmax}

Averaging four encoders instead of three reaches $0.6760$, which is
$0.0011$ above the configuration we report, with a paired interval of
$[-0.0004, +0.0025]$ that includes zero. We report the three-encoder
configuration because the fourth encoder produces no measurable gain
while accounting for close to half of the encoder cost, and because
dropping it also removes a 24-layer model from what has to be run. This
is a choice made on cost and not on score, and we state it as such
rather than reporting the larger number, since both appear in the log
we release.

Dropping that encoder is safe in a second sense. Table~\ref{tab:loeo}
removes each of the four in turn: all four three-encoder averages beat
the single-encoder arm with intervals that exclude zero, so no single
encoder carries the gain. The largest loss comes from dropping the
24-layer model, and even then $0.0046$ of the gain remains. Filipino,
the language that falls under the encoder average, falls less without
that encoder, which is the second reason to leave it out.

% ---------------------------------------------------------------------
% Tabelul `tab:loeo`. Leave-one-encoder-out, din output/e_m3.md.
% ATENTIE la context: familia asta a fost masurata cu ponderea potrivita pe
% splitul de evaluare dar cu modelul lui B inca ales pe trial, deci
% referinta e 0.664581, nu 0.674946. Se scrie in caption, altfel cifrele
% par sa nu se lege de Tabelul 2.
% Perechea patru-vs-trei care conteaza pentru configuratia raportata e in
% ACELASI context ca ea (0.675996 fata de 0.674946) si e in text, nu aici.
% ---------------------------------------------------------------------
\begin{table}[t]
\centering
\small
\begin{tabular}{lcc}
\toprule
Encoders averaged & avg $r$ & vs.\ single \\
\midrule
XLM-R base alone      & .6646 &  \\
All four              & .6716 & $+.0070$ \\
Without XLM-R base    & .6706 & $+.0060$ \\
Without E5            & .6711 & $+.0065$ \\
Without LaBSE         & .6701 & $+.0055$ \\
Without XLM-R large   & .6692 & $+.0046$ \\
\bottomrule
\end{tabular}
\caption{Leave-one-encoder-out. Every three-encoder average beats the
single encoder with a paired interval that excludes zero, the narrowest
being $[+.0011, +.0079]$ for the average we report. These figures were
measured with the fusion weight fitted on the evaluation split but the
language model still chosen on trial, so the reference is $.6646$ rather
than $.6749$.}
\label{tab:loeo}
\end{table}

% =====================================================================
% Sectiunea 5, Conclusions. Scurta, fara idei noi si fara cifre care nu
% sunt deja in Results.
% =====================================================================
\section{Conclusions}
\label{sec:conclusions}

Lexical complexity can be predicted in a language with no annotated
examples at a level that until now required them. Averaged over ten
languages, fusing a frozen transfer regressor with a locally run
language model reaches a Pearson correlation of $0.675$, ahead by
$0.065$ of the strongest published result obtained under the same
zero-label condition, and ahead on nine of the ten languages.

The result rests on a decision procedure rather than on a component. The
same leave-one-language-out criterion that fixes the fusion weight also
decides whether to select a component per language or to average over
all candidates, and it answers the two cases differently: it selects the
language model, where the criterion separates candidates widely, and it
averages the encoders, where the criterion is ten times flatter than the
reference and its maximum is noise. A rule that gave the same answer in
both places would not have been telling us anything about either.

% =====================================================================
% Limitations. Nenumerotata, ca in template. Scrisa imediat dupa Results,
% fiindca e la fel de determinata de masuratori.
%
% Sursele: pasajele 6, 7 si 8 din Paper/TEXT_DE_LIPIT_multilingv.tex, cu
% formularea de bifurcatii venita din output/e_adoptie.md §4. Ultimele doua
% paragrafe sunt ADAUGATE aici, nu erau in pasaje:
%   - trackul combinat, din output/repere_publicate.md, „Regula unica"
%   - cele doua locuri unde LOLO nu e garantat, din
%     output/e_fusion_lolo_adoptat.md, sectiunea de audit
%
% Cifra e 51 de CONFIGURATII (13 pe n curent + 38 in arhiva), nu 70.
% „70" apare inca la finalul lui output/e_m1.md si e numar de randuri de
% jurnal dintr-o stare mai veche. Vezi output/e_adoptie.md §4.
%
% P10 S-A RULAT pe 19 august 2026: `qwen35-9b-8bit` are acum toate cele
% zece limbi pe test, deci bazinul e complet si paragraful despre el a fost
% INLOCUIT, nu sters. Harta ramane Gemma peste tot si cifra de sistem nu s-a
% miscat, deci nimic altceva din lucrare nu se schimba. Marjele si oracolul
% pe catalana vin din output/e_bakeoff_test.md.
% =====================================================================
\section*{Limitations}

We measured 51 system configurations on the test split, listed in full
in the supplementary material. Each of them respects the
leave-one-language-out protocol at every level of choice, so no
per-language figure uses labels from the language it is evaluated on.
The choice of which configuration to report, however, was made by us,
by looking at the mean test score, and is not itself the output of an
automatic procedure. The figure we report should be read as the outcome
of a search of that size rather than a first attempt, in the sense of
\citet{gorman-bedrick-2019-need}. Confidence intervals are paired
against the preceding configuration and are not corrected for multiple
comparisons.

We reduced this exposure where it was cheapest to do so. The single
decision most open to the objection, whether to average encoders or to
select one per language, is now taken by the same criterion used
everywhere else, as described in Section~\ref{sec:selectavg}. The
remaining choices are declared here instead.

All three candidate language models were run on all ten languages of
the test split, so for each language the choice is made over the full
pool rather than over whichever models happened to have been run on it.
The criterion selects Gemma~4 for all ten. Since the criterion for a
language is the mean correlation over the other nine, each of the ten
choices carries a margin over the runner-up, and those margins run from
$0.045$ to $0.060$, so none of them is close. A per-language oracle,
which the protocol does not permit, would differ on one language: on
Catalan, Qwen3.5 reaches $0.507$ where Gemma~4 reaches $0.483$. The
criterion cannot see that, because it never reads the labels of the
language it is choosing for, and the gap is the price of the protocol
rather than a flaw in it.

We evaluate on 5{,}624 of the 5{,}627 test instances. The three
remaining items are cases where the target string could not be located
in its context by any of the matching rules we apply.

We cannot report a figure on the pooled track, in which a single
correlation is computed over all languages at once. Both of our signals
are standardised within the language before they are fused, which is
what lets a single fusion weight serve ten languages, and it leaves the
predictions without a common absolute scale. This is a consequence of
the design rather than a result, but it does mean one published number
has no counterpart here.

Two dependencies survive the protocol, and both are indirect. The
frozen-arm prediction for a fitting language comes from a model trained
on every language except that one, which includes the held-out language,
so the choices made for the held-out language depend on its labels
through predictions on other languages. Measuring the size of that
effect would require training the frozen arm with two languages excluded
at a time, which is ninety fits instead of ten, and we have not run it.
Separately, the language model chosen for a fitting language excludes
one language rather than two. Neither dependency lets a label from a
language enter its own prediction, and the English-only transfer arm has
neither, since it never sees this data at all.

% Custom bibliography entries only
\bibliography{custom}

% Fara apendice la review. Daca se adauga, formatting.md cere sa vina DUPA
% References, cu titlu informativ, iar textul principal sa ramana autonom
% fara ele.

\end{document}
